% Shelved 2026-09-29: continuous-label contrastive pretraining of the FE block.
% Not included in main.tex. The paper is written as if the FE and PE blocks are
% trained without pretraining, until a pretrain-vs-no-pretrain comparison shows
% whether pretraining helps (see kl-nn/repo_reports/kl-paper-gaps.md).
% To restore: paste this subsection back into nn-arch.tex after Sec. 4.1, and
% restore the pretraining sentence at the end of Sec. 4.1.

\subsection{Pretraining the Feature Extractor}
\label{ssec:pretrain}
Before training the PE block, we pretrain the FE block on its own using contrastive learning. Contrastive learning trains a network to map similar inputs to nearby points in an embedding space, and dissimilar inputs to distant points \citep{SimCLR, SupCon}. In our case, we consider two galaxies to be similar if their KL parameters are similar, so that the FE block learns to give similar features to galaxies with similar parameters. Since our parameters are continuous rather than discrete class labels, we use a continuous version of the supervised contrastive loss of \citep{SupCon}, in which each pair of galaxies is weighted by how close their parameters are. \XW{Should state why we pretrain, and whether there is evidence that it helps compared to training the full model from scratch.}

During pretraining, we attach a projection head to the output of the FE block, which is an MLP with layer sizes $1152 \to 2048 \to 512 \to 128$. Its output is normalized to unit length, giving us an embedding $\mathbf{u}_i$ for each galaxy $i$ in a batch. For two galaxies $i$ and $j$, we measure the distance between their parameters as
\begin{equation}
\label{eqn:ccl-distance}
    d_{ij}^2 = \frac{1}{9}\sum_{k=1}^{9}\left(\tilde{\theta}_{ik} - \tilde{\theta}_{jk}\right)^2,
\end{equation}
where $\tilde{\theta}_{ik}$ is the $k$-th parameter of galaxy $i$, rescaled to $[-1, 1]$ using the ranges in Table~\ref{tab:kl-params}. For $\theta_\mathrm{int}$, we take the difference around the circle, so that position angles close to $\pm\pi$ are also treated as close. We then assign a weight to each pair and normalize the weights into a target distribution,
\begin{equation}
\label{eqn:ccl-target}
    w_{ij} = \exp\left(-\frac{d_{ij}^2}{2\sigma_d^2}\right), \qquad
    t_{ij} = \frac{w_{ij}}{\sum_{k\neq i} w_{ik} + w_\mathrm{bg}},
\end{equation}
where $\sigma_d = 0.15$ sets how close two galaxies need to be to count as similar. The background term $w_\mathrm{bg} = \exp\left(-d_\mathrm{cut}^2/2\sigma_d^2\right)$, with $d_\mathrm{cut}=0.4$, makes sure that a galaxy with no other galaxies nearby in parameter space does not put much weight on any of them. Our loss function is then
\begin{equation}
\label{eqn:ccl-loss}
    \mathcal{L}_\mathrm{CCL} = -\frac{1}{N}\sum_{i=1}^{N}\sum_{j\neq i} t_{ij}\log\frac{\exp(\mathbf{u}_i\cdot\mathbf{u}_j/T)}{\sum_{k\neq i}\exp(\mathbf{u}_i\cdot\mathbf{u}_k/T)},
\end{equation}
where $T=0.1$ is a temperature parameter and $N$ is the number of galaxies in the batch. This loss is small when the embeddings of galaxies with similar parameters point in similar directions, and large otherwise. Each galaxy appears in its batch twice, once as simulated and once rotated by $90\degree$, with its parameters transformed accordingly (Sec.~\ref{ssec:npe}). After pretraining, we discard the projection head and keep the weights of the FE block fixed while training the PE block.

